\centerline{\bf \Huge Домашнее задание.}
\centerline{\bf \Huge Комбинаторика. Теория.}

\begin{flushright}
        \scriptsize{Мы становимся сильнее, накапливая результаты.\\ Горечь поражений и наслаждение победами ведут нас к росту.\\ Важно само наличие этих результатов.\\ (c)Аой Тодо. Магическая битва}
\end{flushright}

\problem{1} Известно, что в лососях обучается K различных учеников. У Начальника есть различные олимпиадные задачи их N штук. Сколькими способами можно раздать ученикам все задачки, если:\\
а) K=10, N=20 и кто-то может не получить ни одной задачи?\\
б) K=13, N=12 и каждый должен получить хотя бы одну задачу?\\
в) K=17, N=17 и каждый должен получить хотя бы одну задачу?\\
г) K=4, N=10 и хотя бы один ученик должен быть без задач?\\
д) K=4, N=10 и каждый должен получить хотя бы одну задачу?

\problem{2} У Начальника есть K различных кошельков и N одинаковых купюр. Сколькими способами можно разложить все купюры по кошелькам, если:\\
а) K=10, N=20 и кошельки не могут быть пустыми?\\
б) K=11, N=21 и кошельки могут быть пустыми?